\chapter{Базисно-независимый допуск универсального инцидентного родителя}
\label{ch:v7-universal-incidence-parent-admissibility-gate}

\section{Проверяемая предпосылка предыдущего гейта}

Модулярно-графовый кандидат использовал максимальную бинарную смежность
$A_{\max}$ и получил на копийной плоскости
\begin{equation}
 qA_{\max}^2q
 =3\begin{pmatrix}1&1\\1&1\end{pmatrix}.
 \label{eq:v7-incidence-labeled-square}
\end{equation}
Теперь требуется проверить, является ли эта матрица каноническим следствием
представления и условия первого порядка или уже содержит выбор когерентного
базиса стрелок.

\section{Истинная симметрия изотипической кратности}

Для двух одинаковых бимодулей физическая алгебра действует как
\begin{equation}
 \pi(a)=\pi_0(a)\otimes I_W,
 \qquad W\simeq\mathbb C^2.
 \label{eq:v7-incidence-isotypic-action}
\end{equation}
До задания дополнительного оператора все замены базиса $U\in U(W)$ оставляют
представленные данные неизменными. Поэтому любой канонически и функториально
построенный оператор $T$ на $W$ обязан удовлетворять
\begin{equation}
 UTU^\dagger=T\quad\text{для всех }U\in U(2),
 \qquad T=\lambda I_2.
 \label{eq:v7-incidence-naturality-schur}
\end{equation}
Real-условие может сузить группу до $O(2)$, но полный ортогональный
коммутант на комплексной двумерной плоскости также скалярен.

\section{Независимые стрелки и скрытая когерентная сумма}

Для каждого общего соседа условие первого порядка разрешает две независимые
стрелки, обозначенные ортонормированными векторами $e_1,e_2$ пространства
копий. Базисно-независимая ковариация полного разрешённого пространства
равна
\begin{equation}
 C_{\rm iso}^{(1)}
 =|e_1\rangle\langle e_1|+|e_2\rangle\langle e_2|
 =I_2.
 \label{eq:v7-incidence-incoherent-one-neighbour}
\end{equation}
После трёх общих соседей получается
\begin{equation}
 C_{\rm iso}=3I_2,
 \qquad
 \frac13C_{\rm iso}-I_2=0.
 \label{eq:v7-incidence-incoherent-three-neighbours}
\end{equation}
Это единственный результат, не зависящий от базиса стрелочного пространства.

Напротив, бинарная матрица предыдущего гейта использует для каждого соседа
когерентный столбец $b=e_1+e_2$:
\begin{equation}
 bb^\dagger
 =\begin{pmatrix}1&1\\1&1\end{pmatrix}
 =I_2+\sigma_x.
 \label{eq:v7-incidence-coherent-column}
\end{equation}
Именно перекрёстные члены
$|e_1\rangle\langle e_2|+|e_2\rangle\langle e_1|$ создают
$\sigma_x$. При независимой смене относительной фазы
\begin{equation}
 e_2\longmapsto e^{i\varphi}e_2
 \quad\Longrightarrow\quad
 bb^\dagger-I_2
 =\begin{pmatrix}0&e^{-i\varphi}\\e^{i\varphi}&0\end{pmatrix},
 \label{eq:v7-incidence-relative-phase-orbit}
\end{equation}
поэтому выбор $\varphi=0$ не следует из бинарного факта существования двух
стрелок.

\section{Точный твирлинг}

Усреднение по четырём матрицам Паули уже даёт полный унитарный первый момент:
\begin{equation}
 \frac14\sum_{P\in\{I,\sigma_x,\sigma_y,\sigma_z\}}
 PTP^\dagger=\frac{\Tr T}{2}I_2.
 \label{eq:v7-incidence-pauli-twirl}
\end{equation}
Следовательно,
\begin{equation}
 \operatorname{Twirl}(\sigma_x)=0,
 \qquad
 \operatorname{Twirl}(P_-)=\frac12I_2.
 \label{eq:v7-incidence-selector-twirl}
\end{equation}
Машинный аудит получает нулевые остатки. Усреднение по диэдральному
подмножеству $O(2)$ даёт тот же результат, поэтому Real-структура не
восстанавливает направление.

\section{Граница универсальной дифференциальной оболочки}

Универсальная дифференциальная алгебра $\Omega_{\rm univ}(\mathcal A)$
канонична как алгебраический объект, но её представленная форма использует
коммутатор с уже заданным оператором Дирака:
\begin{equation}
 \pi_D(a_0\,da_1\cdots da_n)
 =\pi(a_0)[D,\pi(a_1)]\cdots[D,\pi(a_n)].
 \label{eq:v7-incidence-represented-universal-form}
\end{equation}
Поэтому она ограничивает допустимые блоки после выбора $D$, но не создаёт
их общую фазу или единичный инцидентный столбец до выбора $D$. Диаграмма
Краевского также кодирует кратности и стрелки конкретной спектральной тройки;
множество всех разрешённых условием первого порядка стрелок не является
каноническим ненулевым оператором Дирака.

\begin{theorem}[Запрет внутреннего вывода единичной инцидентности]
Из представления с кратностью $W\simeq\mathbb C^2$, Real-структуры,
градуировки и линейного пространства всех первопорядково допустимых блоков
канонически получается только скалярная ковариация на $W$. Оператор
$\sigma_x$ предыдущего гейта возникает лишь после когерентного отождествления
двух независимых стрелок, то есть после выбора относительной фазы и
ранга-один в стрелочном пространстве.
\end{theorem}

\begin{nogocriterion}[Запрет скрытого единичного столбца]
Нельзя считать бинарное наличие двух рёбер доказательством того, что их
физические амплитуды равны и фазово синхронизированы. Для открытия маршрута
необходимо отдельное поле или конденсат когерентности стрелок, чьё действие
выводит ранга-один ковариацию и проходит полный гессиан.
\end{nogocriterion}

\begin{protocol}
\begin{enumerate}
 \item полный $U(2)$ изотипической кратности: \textbf{сохраняется};
 \item базисно-независимая ковариация трёх соседей: \textbf{$3I_2$};
 \item нецентральный обмен из пространства всех стрелок: \textbf{нет};
 \item бинарный $A_{\max}$ требует когерентной суммы: \textbf{да};
 \item относительная фаза выведена: \textbf{нет};
 \item Real/O(2) выбирает направление: \textbf{нет};
 \item универсальная оболочка создаёт $D$: \textbf{нет};
 \item прежний условный модулярный расчёт: \textbf{математически сохранён};
 \item его внутренний родительский допуск: \textbf{закрыт};
 \item следующий гейт: \textbf{ранга-один конденсат когерентности стрелок}.
\end{enumerate}
\end{protocol}

Машинный сертификат сохранён в
\path{s2t_v7_universal_incidence_parent_admissibility_gate_results.json}.

\begin{thebibliography}{9}
\bibitem{v7incidence-krajewski}
T.~Krajewski,
\emph{Classification of finite spectral triples},
Journal of Geometry and Physics 28 (1998), 1--30;
arXiv:hep-th/9701081.
\bibitem{v7incidence-cacic}
B.~\v Caci\'c,
\emph{Moduli spaces of Dirac operators for finite spectral triples},
Journal of Noncommutative Geometry 4 (2010), 507--545;
arXiv:0902.2068.
\end{thebibliography}